\section{The Ashkin--Teller reading and open questions}\label{sec:discussion}

Label each direction by the pair of signs $(\sigma,\varsigma)$ of its step in the
diagonal coordinates $U$ and $V$ of Section~\ref{sec:product}, so that east,
north, west and south are $(+,+)$, $(+,-)$, $(-,-)$ and $(-,+)$. Give a
sequence of directions the weight
\[
 \exp\sum_{i<N}\bigl[K(\sigma_i\sigma_{i+1}+\varsigma_i\varsigma_{i+1})
 +K_4\,\sigma_i\sigma_{i+1}\varsigma_i\varsigma_{i+1}\bigr],
\]
the zero-field Ashkin--Teller chain~\cite{ashkinteller1943} with free ends. A
bond that keeps the direction has weight $e^{2K+K_4}$, a turn (one sign flips)
has $e^{-K_4}$, and a reversal (both flip) has $e^{-2K+K_4}$. So the planar walk
is this chain, with a uniform first direction and
\[
 v=e^{-2K-2K_4},\qquad u=e^{-4K}.
\]
Its two magnetizations $\sum_i\sigma_i$ and $\sum_i\varsigma_i$ are $U$ and $V$.
The corner has $U=V=N$, full order of both spins, and the origin has $U=V=0$. So
the crossing compares full order with zero magnetization, as the Ising crossing
of~\cite{paperA} does. The product curve $u=v^2$ is $K_4=0$, where the chain is
two independent Ising chains with the same coupling $K$, and
Proposition~\ref{prop:product} is this classical decoupling read on the walk.
With fixed rates both odds are of order $1/N$, so $K_4=\frac14\log(u/v^2)$
grows like $\frac14\log N$ when $s>0$. The four-spin coupling is then large,
and the regimes of Section~\ref{sec:planar} are those in which it is.

\paragraph{Open questions.}
The survival-and-return estimate of the introduction is not a theorem for
general chains, although we prove its conclusion here. The errors in
Theorem~\ref{thm:fourdir} and Corollary~\ref{cor:planar-turn} are only $o(1)$,
and Theorem~\ref{thm:fourdir} needs $T\le N^{1/8}$. The next term of the
crossing equation when turns dominate is open, and so is the crossover when
$s-2r$ tends to $0$ as $N$ grows. A heuristic suggests that the next term comes
from replacing $T$ by $T+\frac1{4r}-\frac1{2(r+s)}$ in
Theorem~\ref{thm:fourdir}. For $s=0$ this agrees with the term $1/(8z_N)$ of the
crossing equation in Section~\ref{sec:notation}, through
Theorem~\ref{thm:pureturns}(c). Off the product curve we have no closed form for
$\omega_N$ or for the crossing, and Corollary~\ref{cor:productcompare} only brackets
the crossing. With turns only, the computations give $N(v^*_N-u_N)=0.37\,L^2/N$
at $N=96$ and $0.38\,L^2/N$ at $N=10240$. So the order $L^2/N$ in
Theorem~\ref{thm:pureturns}(c) seems to be right, but we do not prove a lower
bound of this order. The Ashkin--Teller chain at fixed couplings $K,K_4$ with
$N\to\infty$, the clock and Potts chains with more states, and the walk on
$\Z^d$ with $d\ge3$ are open as well.
